\documentclass[10pt,a4paper]{amsart}
\usepackage[margin=19mm]{geometry}
\usepackage{amsmath,amssymb,mathtools}
\usepackage[scr=rsfs,cal=euler]{mathalfa}
\usepackage{xfrac}
\usepackage[unicode,hidelinks,bookmarks=false]{hyperref}
\newcommand{\ie}{i.e.\ }
\newcommand{\cf}{cf.\ }
\newcommand{\resp}{respectively\ }
\newcommand{\Subsection}{\subsection}
\providecommand{\nobreakdash}{\nobreak\hbox{-}}
\setlength{\emergencystretch}{2em}
\multlinegap0pt
\usepackage{mathtools}
\newtheorem{theorem}{Theorem}
\newcommand{\Comp}{\mathbb{C}}
\newcommand{\Q}{\mathbb{Q}}
\newcommand{\Z}{\mathbb{Z}}
\newcommand{\image}{\operatorname{Im}}
\newcommand{\res}{\operatorname{Res}}
\title{Geometric construction of modular polynomials with level structures}
\author{Hiroshi Onuki, Yukihiro Uchida, Ryo Yoshizumi}
\date{Source: arXiv:2601.17338v1 (CC BY 4.0)}
\begin{document}
\maketitle

\noindent\emph{Source and licence.} Geometric construction of modular polynomials with level structures. Version arXiv:2601.17338v1 (CC BY 4.0), distributed by arXiv under the Creative Commons Attribution 4.0 International licence, http://creativecommons.org/licenses/by/4.0/, as recorded in the arXiv licence metadata for this exact version and shown on the abstract page https://arxiv.org/abs/2601.17338v1. Full licence notice: \texttt{licenses/arxiv-2601.17338-CC-BY-4.0.txt}.\par\smallskip

\noindent\textit{Editorial scope.} Counted unit: the article's theorem thm:integral\_modular\_polynomial (source0.tex lines 621--627). The article's complete original proof of it is delivered with it (source0.tex lines 629--702, one \texttt{proof} environment of 74 lines). No mathematics is written, repaired or re-derived: the statement and the proof are the article's own lines, in the article's own notation, and no retained line is substituted or re-flowed.  Retained uncounted as the source-local context the proof needs: lines 592--594.  Nothing was omitted from any retained span, so the package carries no omission record.  No retained line contains a citation, so the wrapper carries no bibliography and no bibliography entry.  No retained line contains a figure environment, an image-inclusion command or drawing code, so no image is delivered and the package declares no asset dependency.  Editorial formatting only, in addition: an AMS \texttt{amsart} preamble that restates the article's own declarations and macros used by the retained lines, namely the environments \texttt{theorem} on separate counters, together with the standard packages those lines need.  The retained environments print in this document as Theorem 1 in order of appearance, whereas the article prints the same items with the numbering of its own full document.  The complete native source of the article as distributed in the arXiv e-print for this exact version is bundled unchanged as \texttt{sources/193308/source0.tex}.
    \begin{equation}\label{eq:def_mod_poly}
        F(Y) = \prod_{i=1}^{\psi(m)} (Y - \sigma_i),
    \end{equation}

\begin{theorem}\label{thm:integral_modular_polynomial}
    The modular polynomial $\Phi_m^{\alpha}(X, Y)$ satisfies the following:
    \begin{enumerate}
        \item $\Phi_m^{\alpha}(X, Y) \in \Z[X, Y]$. \label{enum:integer_coeff}
        \item $\deg_X(\Phi_m^{\alpha}) = \psi(m)$. \label{enum:degree_X}
    \end{enumerate}
\end{theorem}

\begin{proof}
    (\ref{enum:integer_coeff}):
    Recall the definition of $\Phi_m^{\alpha}(X, Y)$ in \eqref{eq:def_mod_poly}.
    Since the leading coefficient of $\Psi_m(X; Y)$ in $Y$ is an integer not divisible by any prime not dividing $m$,
    and $V_m$ has coefficients in $\Z\left[\frac{1}{m}\right]$,
    the roots $\sigma_1, \dots, \sigma_{\psi(m)}$ of $\Phi_m^{\alpha}(X, Y)$
    as a polynomial in $Y$
    are integral over $\Z\left[\frac{1}{m}, X\right]$.

    Let $j_1, \dots, j_{\psi(m)}$ be the roots of
    the classical modular polynomial $\Phi_m\left(\frac{J_1(X)}{J_0(X)}, Y\right)$
    in an algebraic closure of $\Q(X)$.
    Then, we have $j_i = J_1(\sigma_i)/J_0(\sigma_i)$ for all $1 \leq i \leq \psi(m)$
    by reindexing if necessary.
    Since $\Phi_m$ has integer coefficients and is monic in $Y$,
    the root $j_i$ is integral over $\Z\left[\frac{J_1(X)}{J_0(X)}\right]$ for all $1 \leq i \leq \psi(m)$.
    Let $c$ be the leading coefficient of $J_1(X)$.
    Then, $\sigma_i$ is integral over $\Z\left[\frac{1}{c}, j_i\right]$
    since we have $J_1(\sigma_i) - J_0(\sigma_i) j_i = 0$.
    Thus, 
    $\sigma_i$ is integral over $\Z\left[\frac{1}{c}, \frac{J_1(X)}{J_0(X)}\right]$.

    Combining the above results,
    we see that $\sigma_i$ is integral over $\Z\left[\frac{1}{m}, X\right] \cap \Z\left[\frac{1}{c}, \frac{J_1(X)}{J_0(X)}\right]$.
    Since $m$ is prime to $c$ and $J_0(X)$ is primitive,
    we have $\Z\left[\frac{1}{m}, X\right] \cap \Z\left[\frac{1}{c}, \frac{J_1(X)}{J_0(X)}\right] = \Z[X]$.
    This implies that the coefficients of $\Phi_m^{\alpha}(X, Y)$ are in $\Z$.

    (\ref{enum:degree_X}):
    Let $F(X, Y) = J_1(X) - J_0(X) Y$,
    and $d = \deg J_1$.
    Consider four variables $X, Y, Z, W$,
    and define a polynomial $G(X, Y) \in \Z[X, Y]$ by
    \begin{equation*}
        G(X, Y) = \res_Z(\res_W(F(W, X), \Phi_m^\alpha(W, Z)), F(Z, Y)) \in \Z[X, Y].
    \end{equation*}
    Here, $\res_Z$ and $\res_W$ denote the resultant with respect to
    variables $Z$ and $W$, respectively.
    By the property of resultant,
    for any $j, j' \in \Comp$,
    it holds that $G(j, j') = 0$
    if and only if there exist $z, z' \in \Comp$
    such that $F(z, j) = 0$, $F(z', j') = 0$,
    and $\Phi_m^\alpha(z, z') = 0$.
    If $z \not\in \image \alpha$ then $J_0(z) = 0$,
    and thus $F(z, j) \neq 0$ for any $j \in \Comp$.
    This means that $G(j, j') = 0$ if and only if
    there exist $z, z' \in \image \alpha$
    such that $E_z^\alpha$ and $E_{z'}^\alpha$ are connected by
    a cyclic isogeny of degree $m$,
    and $j = j(E_z^\alpha)$, $j' = j(E_{z'}^\alpha)$.
    Therefore, $G(j, j') = 0$ if and only if $\Phi_m(j, j') = 0$.
    From the irreducibility of $\Phi_m(X, Y)$,
    it follows that $G(X, Y) = \gamma \cdot \Phi_m(X, Y)^n$ for some $\gamma \in \Comp^\times$ and $n \in \Z_{\geq 1}$.

    For $j \in \Comp$, we let $\omega_1, \dots, \omega_d$
    be the roots of $F(W, j)$ in $\Comp$.
    For $i = 1, \dots, d$,
    let $z_{i, 1}, \dots, z_{i, \psi(m)}$ be the roots of $\Phi_m^\alpha(\omega_i, Z)$ in $\Comp$.
    Then, we have
    \begin{align*}
        G(j, Y) &= \res_Z\left(c^{\deg_X(\Phi_m^\alpha)}\prod_{i=1}^{d} \Phi_m^\alpha(\omega_i, Z), F(Z, Y)\right)\\
            &= c^{d\cdot\deg_X(\Phi_m^\alpha)} \prod_{i=1}^{d} \prod_{j=1}^{\psi(m)} F(z_{i, j}, Y).
    \end{align*}
    Therefore, we have $\deg_Y(G(j, Y)) = d \cdot \psi(m)$.
    Since this holds for any $j \in \Comp$,
    we have $\deg_Y(G(X, Y)) = d \cdot \psi(m)$.
    Thus, we have $G(X, Y) = \gamma \cdot \Phi_m(X, Y)^d$.

    Similar discussion considering $G(X, j)$ for $j \in \Comp$
    shows that $\deg_X(G(X, Y)) = d\cdot\deg_X(\Phi_m^\alpha)$.
    By comparing the degrees of $G(X, Y)$ and $\Phi_m(X, Y)^d$ in $X$,
    we obtain $\deg_X(\Phi_m^\alpha) = \psi(m)$.
\end{proof}

\end{document}
